\documentclass[aspectratio=1610,xcolor=dvipsnames]{beamer}

\usepackage{etex}
\usepackage[utf8]{inputenc}
\usepackage{tgtermes}
\usefonttheme{default}
\renewcommand{\familydefault}{\rmdefault}
\usetheme{Madrid}
\usepackage{lmodern}

% --- Ethereum brand palette (blue layout + purple emphasis) ---
\definecolor{ethpurple}{HTML}{627EEA}   % layout blue (headers, footline, bullets)
\definecolor{ethdeep}{HTML}{454A75}      % deep indigo-purple (structure contrast)
\definecolor{ethblue}{HTML}{88AAF1}      % logo light blue (secondary accent)
\definecolor{ethdark}{HTML}{3C3C3D}     % Ethereum charcoal (neutral text)
\definecolor{ethaccent}{HTML}{7B3FE4}   % emphasis purple (ETH gradient violet)
\definecolor{canvasbg}{HTML}{ECEFF0}     % Ethereum light gray background
\definecolor{graphite}{HTML}{3C3C3D}     % body text
\colorlet{ethshade}{ethpurple!80!black}
\colorlet{canvasdark}{canvasbg!80!black}

% --- Beamer colors ---
\setbeamercolor{structure}{fg=ethpurple}
\setbeamercolor{background canvas}{bg=canvasbg}
\setbeamercolor{frametitle}{bg=ethpurple, fg=white}
\setbeamercolor{title}{bg=ethpurple, fg=white}
\setbeamercolor{subtitle}{fg=white,bg=ethpurple}
\setbeamercolor{item}{fg=ethpurple}
\setbeamercolor{section in toc}{fg=ethdeep}
\setbeamercolor{subsection in toc}{fg=ethpurple}
\setbeamercolor{author in head/foot}{bg=ethpurple, fg=white}

% Footline (Madrid dark bar): white text for short title and page number
\setbeamercolor{title in head/foot}{fg=white}
\setbeamercolor{date in head/foot}{fg=white}


\usepackage{tikz}
\usetikzlibrary{shapes.geometric, arrows.meta, positioning}
\usepackage{graphicx} % For \resizebox

% --- TikZ styles (Ethereum palette) ---
\tikzset{
  block/.style = {rectangle, draw, fill=ethpurple!40, text width=4cm, align=center, rounded corners, minimum height=2.0em},
  sideblock/.style = {rectangle, draw, fill=ethblue!40, text width=4cm, align=center, rounded corners, minimum height=2.0em},
  decision/.style = {diamond, draw, fill=canvasbg!40, text width=3cm, align=center, aspect=2, inner sep=1pt},
  arrow/.style = {thick, ->, >=Stealth},
  mechanismblock/.style = {rectangle, draw, fill=ethblue!20, text width=4cm, align=center, minimum height=2.5em},
  greydashedarrow/.style = {thick, dashed, ->, >=Stealth, draw=gray},
  blusharrow/.style = {thick, ->, >=Latex, draw=ethdeep},
  blank/.style = {rectangle, fill=ethdeep!20, text width=4cm, align=center, minimum height=2.0em}
}

%<---- do not use enumitem, does not work well with Beamer

\usepackage{amsmath, amssymb, amsfonts}
\usepackage{booktabs, array, dcolumn}
\usepackage{graphicx, subfig}
\usepackage{epstopdf}
\def\pdfshellescape{1}
\usepackage{ragged2e}
\AtBeginEnvironment{frame}{\justifying}

% Icons/bullets
\setbeamertemplate{itemize items}[ball]
\setbeamertemplate{itemize subitem}[triangle]

% Theorems
\usepackage{amsthm}
\newtheorem{proposition}{Proposition}
\newtheorem*{quiz}{Quiz}
\newtheorem{claim}[proposition]{Claim}
\newtheorem{exercise}[proposition]{Exercise}
\newtheorem{remark}[proposition]{Remark}

\usepackage{hyperref}
\usepackage{appendixnumberbeamer}
\definecolor{links}{HTML}{627EEA}
\hypersetup{colorlinks,linkcolor=links,urlcolor=links,citecolor=links}

\setbeamertemplate{navigation symbols}{}
\newcommand{\footlinemeeting}{3rd WESEAMS}
\newcommand{\backlink}[2]{\hyperlink{#1}{\textcolor{links}{$\leftarrow$ #2}}}
\newcommand{\fwdlink}[2]{\hyperlink{#1}{\textcolor{links}{#2 $\rightarrow$}}}
\setbeamertemplate{footline}{%
  \leavevmode%
  \hbox{%
    \begin{beamercolorbox}[wd=.72\paperwidth,ht=2.25ex,dp=1ex,leftskip=1cm]{title in head/foot}%
      \insertshorttitle \hspace{0.8em}---\hspace{0.8em}\footlinemeeting
    \end{beamercolorbox}%
    \begin{beamercolorbox}[wd=.28\paperwidth,ht=2.25ex,dp=1ex,rightskip=1cm]{date in head/foot}%
      \hfill\insertframenumber\,/\,\inserttotalframenumber
    \end{beamercolorbox}%
  }%
}
\newcommand{\hilight}[1]{\colorbox{ethblue!35}{#1}}

\DeclareMathOperator*{\plim}{plim}
\DeclareMathOperator*{\argmax}{arg\,max}
\DeclareMathOperator*{\E}{E}
\DeclareMathOperator*{\Var}{Var}
\DeclareMathOperator*{\Cov}{Cov}
\DeclareMathOperator*{\Corr}{Corr}
\DeclareMathOperator*{\supp}{supp}


\newcommand{\ind}{\mathrel{\perp \! \! \! \perp}}
\def\citeapos#1{\citeauthor{#1}'s (\citeyear{#1})}

\title[Student Sharing 2026]{Student Sharing}
\subtitle{HKUST Department of Economics\\
PhD Early Admission Camp 2026}
\author[Harlly Zhou]{Harlly Zhou}
\institute[HKUST]{Department of Economics\\
HKUST Business School}
\date{}

\begin{document}

%%%%%%%%%%%%%%%%%%%%%%%%%%%%%%%%%%%%%%%%%%%%
%%%%%%%%%%%%%%%%%%%%%%%%%%%%%%%%%%%%%%%%%%%%
\begin{frame}
\titlepage
\end{frame}

\begin{frame}{About Myself}
    \begin{itemize}
        \item Undergraduate at HKU
        \item Master at UChicago
        \item Supervisor: Prof. Yang Lu
        \item Research Interests: (Macroeconomics) Monetary Policy, Production Network
    \end{itemize}
\end{frame}

\begin{frame}{Why Hong Kong?}
    \begin{itemize}
        \item Fully funded: Studentship, HKPFS, other scholarships
        \item Permanent residence: 7 years of visa
        \item Active research community: 8 main universities with many interactions and collaborations
    \end{itemize}
\end{frame}

\begin{frame}{Why HKUST?}
    \begin{itemize}
        \item Research university: World-class faculties
        \item Supervision: Flexible and supportive
        \item Community: Collaborations among students and faculties
    \end{itemize}
\end{frame}

\begin{frame}{Really want a PhD?}
    \begin{itemize}
        \item Love research? Or you do not want to work?
        \item STRONG motivation: Research should not be mentally-detrimental to you.
        \item Physical and mental health: Research is a long-term commitment.
    \end{itemize}
\end{frame}

\begin{frame}{What a PhD needs? From a MBTI perspective}
    \begin{itemize}
        \item Introvert vs Extrovert
        
        Introvert for doing research along; Extrovert for presentation, socializing, and collaboration.

        \item Sensing vs Intuition
        
        Sensing for doing research step by step; Intuition for innovative ideas.

        \item Thinking vs Feeling
        
        Thinking for logical reasoning; Feeling for emotion management.

        \item Judging vs Perceiving
        
        Judging for planning and organization; Perceiving for flexibility and adaptability.
    \end{itemize}
\end{frame}


\end{document}